% main.tex
% An Executable Reproduction Artifact for Equation (2.2)
% of the Erdős Unit-Distance Disproof Remarks Paper
%
% Compile:  latexmk -pdf main.tex   (uses biber for the bibliography)
%
% Style: deliberately close to the typographic conventions of
% unit-distance-proof.pdf (OpenAI 2026), with the document TYPE marked
% as a reproduction artifact rather than original research. Theorem
% environments are replaced by "Reproduction" boxes; the paper does not
% introduce any new mathematical result.

\documentclass[11pt,a4paper]{article}

% --- Packages ---------------------------------------------------------------
\usepackage[utf8]{inputenc}
\usepackage[T1]{fontenc}
\usepackage{lmodern}
\usepackage[margin=1in]{geometry}
\usepackage{amsmath, amssymb, amsthm}
\usepackage{microtype}
\usepackage{enumitem}
\usepackage{hyperref}
\usepackage{xcolor}
\usepackage{booktabs}
\usepackage[backend=biber,style=alphabetic,sorting=nyt]{biblatex}

\addbibresource{references.bib}

\hypersetup{
  colorlinks=true,
  linkcolor=black,
  citecolor=black,
  urlcolor=blue!50!black,
  pdftitle={An Executable Reproduction Artifact for Equation (2.2) of the Erdős Unit-Distance Disproof Remarks Paper},
  pdfauthor={Yun Kwansub}
}

% --- Reproduction environment (replaces \newtheorem{theorem}) ---------------
% We deliberately do NOT define a Theorem environment in this paper:
% nothing in this artifact constitutes a new theorem. The Reproduction
% environment records a specific numerical or combinatorial value that
% this artifact computes and pins against the published source.
\newtheoremstyle{reproduction}%
  {1.2em}{1.2em}%        % top/bottom spacing
  {\itshape}{}%          % body / indent
  {\bfseries}{.}%        % heading / punctuation
  {.5em}%                % space after head
  {}                     % default head spec
\theoremstyle{reproduction}
\newtheorem{reproduction}{Reproduction}[section]
% Backward-compatible alias so existing \begin{verification} blocks render.
\let\verification\reproduction
\let\endverification\endreproduction

% Macros
% Breakable monospace identifier (long test / function / path names with
% underscores). Built on \DeclareUrlCommand from the url package (already
% loaded transitively by hyperref): the argument is taken verbatim, so
% underscores need no backslash-escape, and line breaks are permitted
% after _, /, ., :, -, ( and ).  Write \tname{foo_bar}, not \tname{foo\_bar}.
\DeclareUrlCommand{\tname}{\urlstyle{tt}%
  \def\UrlBreaks{\do\.\do\/\do\_\do\-\do\:\do\(\do\)}}
\newcommand{\Q}{\mathbb{Q}}
\newcommand{\Z}{\mathbb{Z}}
\newcommand{\C}{\mathbb{C}}
\newcommand{\R}{\mathbb{R}}
\newcommand{\OK}{\mathcal{O}_K}
\DeclareMathOperator{\Gal}{Gal}
\DeclareMathOperator{\Disc}{Disc}

% --- Title block ------------------------------------------------------------
\title{%
  An Executable Reproduction Artifact for\\
  Equation (2.2) of the Erd\H{o}s Unit-Distance\\
  Disproof Remarks Paper}
\author{%
  Yun Kwansub%
  \thanks{Independent author. Code distribution:
    \href{https://github.com/Flamehaven-Labs/erdos-ant-verification}{github.com/Flamehaven-Labs/erdos-ant-verification}.}%
}
\date{Version 0.2.0 --- 2026-05-25}

% --- Document --------------------------------------------------------------
\begin{document}
\maketitle

\begin{abstract}
We present \texttt{erdos-ant-verification}, an open-source Python package
that re-derives, in pure NumPy and \texttt{mpmath}, the finite and
numerically explicit parts of the 2026 disproof of the Erd\H{o}s planar
unit-distance conjecture announced by an OpenAI reasoning model
\cite{OpenAI2026Proof} and written up by Alon, Bloom, Gowers, Litt,
Sawin, Shankar, Tsimerman, Wang, and Matchett Wood
\cite{Sawin2026Remarks}. The central object is equation~(2.2) of the
remarks paper: for the explicit choice
$T = \{3, 5, 7, 11, 13, 17\}$, $S = \{101, \infty\}$, and
$L_T = \Q(\sqrt{5}, \sqrt{13}, \sqrt{17}, \sqrt{21}, \sqrt{33})$,
that equation states a rigorous numerical lower bound
$\delta_{\text{excess}} \geq \approx 6.24 \times 10^{-38}$ on the excess
above the classical Erd\H{o}s exponent. Our \texttt{mpmath} evaluation
at $200$ bits of precision returns $6.2391 \times 10^{-38}$, matching
the published two-significant-figure value to $1.4 \times 10^{-4}$
relative error.

We make no other mathematical claim. The artifact does \emph{not}
re-prove Theorem~1.1 of \cite{Sawin2026Remarks}, does \emph{not}
construct the infinite Golod--Shafarevich tower whose existence is
quoted as an input to that theorem, and does \emph{not} address any
separately published sharpening of the bound (in particular, it does
\emph{not} reproduce or verify the explicit value
$\delta \approx 0.014$ that appears in Sawin's separate explicit-lower-bound preprint or in informal coverage of the result).
The intended audience is anyone who, having read \S2 of
\cite{Sawin2026Remarks}, wants a citable, CI-tested Python
implementation of its eq.~(2.2) on which they can run their own
sensitivity checks. The package is distributed under the MIT licence
at \url{https://github.com/Flamehaven-Labs/erdos-ant-verification}.

\noindent\textbf{Reproduction.}
After cloning the repository,
\verb|pip install -e .[dev]| followed by
\verb|python -m erdos_ant.verify| runs the test suite and prints the
verdict; the same command with \verb|--write-evidence| also refreshes
the tracked JSON certificate at
\path{reports/verification_result.json}. The frozen evidence for this
paper corresponds to commit \texttt{c8fe529} (release tag
\texttt{v0.1.3}) with Zenodo DOI
\href{https://doi.org/10.5281/zenodo.20377950}{10.5281/zenodo.20377950}.\footnote{The earlier source-only release \texttt{v0.1.1} carries DOI \href{https://doi.org/10.5281/zenodo.20366884}{10.5281/zenodo.20366884}. The intermediate \texttt{v0.1.2} was a pre-PDF snapshot, and \texttt{v0.1.3}/\texttt{v0.1.4} bundled the first compiled paper. This document is \texttt{v0.2.0}, a documentation-only revision that tightens scope wording and removes terminology specific to the author's internal tooling; the underlying code, tests, and numerical claim are unchanged.}
\end{abstract}

\clearpage
\tableofcontents
\clearpage

% ---------------------------------------------------------------------------
\section{Background and scope}
\label{sec:background}

% Section 1: Background and scope

\subsection{The Erd\H{o}s unit-distance problem}

For a finite set $P \subset \R^2$, let
$\nu(P) = \#\{\{x, y\} \subset P : |x - y| = 1\}$
and $\nu(n) = \max_{|P| = n} \nu(P)$. The planar unit-distance problem
\cite{Erdos1946} concerns the asymptotic behaviour of $\nu(n)$. Erd\H{o}s
proved $\nu(n) \geq n^{1 + \Omega(1/\log \log n)}$ via a square-grid
construction and conjectured the matching upper bound
$\nu(n) \leq n^{1 + C / \log \log n}$ for some absolute constant $C$.

\subsection{The 2026 disproof}

The conjecture was disproved in May 2026 by an internal reasoning model
at OpenAI \cite{OpenAI2026Proof}, with a simplified construction and
explicit numerical lower bound published the following day by Alon,
Bloom, Gowers, Litt, Sawin, Shankar, Tsimerman, Wang, and Matchett Wood
\cite{Sawin2026Remarks}. The combined result is:

\medskip
\noindent\textbf{Theorem 1.1 of \cite{Sawin2026Remarks}.}
\textit{There exists an absolute constant $\delta > 0$ and infinitely many
positive integers $n$ such that $\nu(n) \geq n^{1 + \delta}$.}

\medskip
The construction passes through a CM field $K_j = L_j(i)$, where the
$L_j$ form an infinite tower of totally real fields of bounded root
discriminant, supplied by the Golod--Shafarevich theorem
\cite{GolodShafarevich1964}. Section~2 of \cite{Sawin2026Remarks}
gives an explicit choice with
$T = \{3, 5, 7, 11, 13, 17\}$, $S = \{101, \infty\}$, and
$L_T = \Q(\sqrt{5}, \sqrt{13}, \sqrt{17}, \sqrt{21}, \sqrt{33})$,
and computes (eq.~(2.2)) a rigorous numerical lower bound
$\delta \geq \approx 6.24 \times 10^{-38}$.

\subsection{What this artifact verifies}
\label{sec:scope}

This paper accompanies a Python package, \texttt{erdos-ant-verification},
whose purpose is narrow and explicitly bounded:

\begin{itemize}[leftmargin=*]
  \item we \emph{do not} re-prove Theorem~1.1;
  \item we \emph{do not} construct the infinite Golod--Shafarevich tower;
  \item we \emph{do} re-implement the finite, explicitly checkable parts
        of the construction in pure Python, with exact small-integer
        arithmetic for the algebraic-number-theory layer and
        \texttt{mpmath} at 200-bit precision for the numerical evaluation
        of eq.~(2.2);
  \item we \emph{do} pin the published numerical value
        $\approx 6.24 \times 10^{-38}$ as a regression test that fails if
        our code drifts from the published source.
\end{itemize}

The artifact is therefore neither a proof checker (it does not formalise
the proof in a trusted kernel) nor an independent mathematical
contribution (it produces no new theorem). It is a reproducibility
artifact in the sense of, e.g., the ACM Artifact Evaluation guidelines:
a citable, version-pinned, CI-verified executable that allows any
third party to re-derive the published numerical value on their own
hardware without relying on our reading of the source PDFs.

The remainder of the paper describes the three implementation phases
(\S\S\ref{sec:phase1}--\ref{sec:phase3}), the central numerical
reproduction (\S\ref{sec:eq22}), the explicit limits
(\S\ref{sec:limits}), and the reproducibility surface
(\S\ref{sec:repro}).


% ---------------------------------------------------------------------------
\section{Phase 1: faithful $h(K) = 1$ lattices}
\label{sec:phase1}

% Section 2: Phase 1 -- faithful h(K)=1 lattices

Phase~1 implements the trivial-class-number degenerate case of Lemma~2.2
of \cite{OpenAI2026Proof}: an imaginary quadratic field $K = \Q(\sqrt{-d})$
with class number $h(K) = 1$, where the ideal-class pigeonhole reduces to
the identity. For these fields the ring of integers $\OK$ embeds
standardly in $\R^2$ as a rank-$2$ lattice, and unit distances between
lattice points can be counted exactly.

\subsection{Coverage}

We implement the three smallest $h = 1$ cases (the Baker--Heegner--Stark
list \(1, 2, 3, 7, 11, 19, 43, 67, 163\) is encoded in the module-level
constant \texttt{CLASS\_NUMBER\_ONE\_DISCRIMINANTS}; the first three are
implemented):

\begin{tabular}{lll}
\toprule
$d$ & $\OK$ & Standard embedding into $\R^2$ \\
\midrule
$1$ & $\Z[i]$ (Gaussian integers)     & $a + bi \mapsto (a, b)$ \\
$2$ & $\Z[\sqrt{-2}]$                 & $a + b\sqrt{-2} \mapsto (a, b\sqrt{2})$ \\
$3$ & $\Z[\omega]$, $\omega = \frac{-1 + \sqrt{-3}}{2}$ & $a + b\omega \mapsto (a - b/2,\; b\sqrt{3}/2)$ \\
\bottomrule
\end{tabular}

\medskip
For each $d$ and bound $b \geq 0$, the module enumerates all
$(2b+1)^2$ lattice points with integer coefficients in $[-b, b]$
and counts unit-distance pairs exactly via a numpy $O(n^2)$
broadcast.

\subsection{Verifications}

\begin{verification}\label{ver:gauss_3x3}
For $\Z[i]$ with $b = 1$ (the $3 \times 3$ Gaussian grid, $n = 9$),
the artifact computes exactly $12$ unit-distance pairs: three rows of two
horizontal edges plus three columns of two vertical edges. Pinned in
\tname{tests/test_imaginary_quadratic_lattice.py::test_gaussian_3x3_grid_has_12_unit_distances}.
\end{verification}

\begin{verification}\label{ver:gauss_7x7}
For $\Z[i]$ with $b = 3$ (the $7 \times 7$ Gaussian grid, $n = 49$),
the artifact computes exactly $2 \cdot 7 \cdot 6 = 84$ unit-distance
pairs.
\textit{Source:
\tname{src/erdos_ant/imaginary_quadratic_lattice.py},
class \tname{ImaginaryQuadraticLattice.count_unit_distance_pairs}.}
Pinned in \tname{src/erdos_ant/verify.py} as the assertion
\tname{phase1_gaussian_7x7_grid_has_expected_84_pairs}.
\end{verification}

\begin{verification}\label{ver:eisenstein_dominance}
For each $b \in \{1, 2, 3, 4\}$, the Eisenstein-lattice (triangular)
unit-distance count strictly exceeds the Gaussian-lattice count at the
same $b$. The triangular lattice has $6$ nearest neighbours per
interior point versus the square lattice's $4$. Pinned in
\tname{test_eisenstein_unit_distances_exceed_gaussian_for_same_bound}.
\end{verification}

\subsection{Non-claims}

Phase~1 does not improve the asymptotic Erd\H{o}s exponent. With
$[K : \Q] = 2$, the construction is the classical square or triangular
grid; in particular $\nu(n) \leq n^{1 + O(1/\log \log n)}$ holds and the
$\delta > 0$ excess is not present. The asymptotic result strictly
requires $[K : \Q] \to \infty$, which is the Phase~3 setting.

Phase~1 is included in the artifact for two reasons: (i) it is a
self-contained, exact reference against which the Phase~3 construction
can be sanity-checked at low parameter values, and (ii) it documents
the standard complex embedding of $\OK$ in $\R^2$ unambiguously, which
is the substrate of every later step.


% ---------------------------------------------------------------------------
\section{Phase 2: $K = \Q(\sqrt{-5})$ with genus theory}
\label{sec:phase2}

% Section 3: Phase 2 -- Q(sqrt(-5)) with class group + genus theory

Phase~2 implements the smallest non-trivial-class-number case
$K = \Q(\sqrt{-5})$, $h(K) = 2$. The Lemma~2.2 pigeonhole of
\cite{OpenAI2026Proof} becomes non-vacuous here: when an ideal product
$A_\varepsilon = \prod_j P_j^{\varepsilon_j} (P_j')^{k - \varepsilon_j}$
lands in the non-principal class, no element $\alpha_\varepsilon \in K^*$
with $(\alpha_\varepsilon) = A_\varepsilon A_\eta^{-1}$ exists for
arbitrary $\eta$ --- the pigeonhole must restrict to a single ideal
class, and the bound $|U| \geq (k+1)^s / h(K)$ is sharp.

\subsection{Classical facts encoded}

\begin{itemize}[leftmargin=*]
  \item $\OK = \Z[\sqrt{-5}]$ with discriminant $-20$
        (\texttt{K\_DISCRIMINANT = -20}, \texttt{CLASS\_NUMBER = 2}).
  \item The non-trivial ideal class is represented by
        $\mathfrak{p}_2 = (2, 1 + \sqrt{-5})$, satisfying
        $\mathfrak{p}_2^2 = (2)$.
  \item By genus theory \cite[Theorem~6.1]{Cox1989}, the Hilbert class
        field is the genus field $H = K(\sqrt{-1}) = \Q(i, \sqrt{5})$.
  \item A rational prime $p$ (odd, $p \neq 5$) splits completely in $K$
        iff $\left(\frac{-5}{p}\right) = 1$, equivalently
        $p \bmod 20 \in \{1, 3, 7, 9\}$. Of these, $p \bmod 20 \in \{1, 9\}$
        gives a principal prime (representable as $a^2 + 5b^2$), and
        $p \bmod 20 \in \{3, 7\}$ gives a non-principal prime.
\end{itemize}

\subsection{Lemma~2.2 pigeonhole, explicit}

The Phase~2 module \texttt{erdos\_ant.genus\_class\_field} exposes a
\texttt{GenusClassFieldProbe} that, given a tuple of split rational
primes, computes the ideal-class partition of the
$\{0,1\}^s$ family of ideal products and reports
$|U| \geq (k+1)^s / h(K)$.

\begin{verification}\label{ver:lemma22_principal}
For the choice $(p_1, p_2) = (29, 41)$ (both in the principal genus,
$29 = 3^2 + 5 \cdot 2^2$, $41 = 6^2 + 5 \cdot 1^2$), all four candidate
ideals $A_\varepsilon$ for $\varepsilon \in \{0,1\}^2$ are principal,
and the Lemma~2.2 lower bound is
$(1 + 1)^2 / 2 = 2$. Pinned in
\tname{tests/test_genus_class_field.py::test_lemma_22_pigeonhole_with_two_principal_primes}.
\end{verification}

\begin{verification}\label{ver:lemma22_mixed}
For the choice $(p_1, p_2) = (29, 3)$ (one principal, one
non-principal), the class group has order $2$, so the prime above $3$
and its conjugate lie in the same non-principal class. Consequently all
four ideals land in the non-principal class rather than splitting
$2 + 2$ between the two classes, and the Lemma~2.2 lower bound is still
$2$.
Pinned in \tname{test_lemma_22_pigeonhole_with_mixed_genus}.
\end{verification}

\subsection{Non-claims}

Phase~2 fixes a single CM field of degree $2$ over $\Q$ and produces a
finite pigeonhole demonstration. It does not produce a sequence of
fields with $[K_j : \Q] \to \infty$, and does not generate the
exponentially growing set $U$ that drives the asymptotic exponent
excess. It does, however, instantiate the smallest concrete example
where the pigeonhole step is genuinely non-trivial.


% ---------------------------------------------------------------------------
\section{Phase 3: the Sawin multi-quadratic construction, finite parts}
\label{sec:phase3}

% Section 4: Phase 3 -- Sawin multi-quadratic construction (finite parts)

Phase~3 implements the finite, explicitly checkable parts of the
\emph{simplified} construction of Section~2.1 of
\cite{Sawin2026Remarks}, which uses an unramified pro-$2$ tower over a
multi-quadratic base field rather than the original GPT pro-$3$ tower
over a cyclic cubic. The simplification is due to Sawin and yields the
explicit numerical bound evaluated in \S\ref{sec:eq22}.

\subsection{The Sawin parameters, verbatim}

The construction fixes
\begin{align*}
  T   &= \{3,\, 5,\, 7,\, 11,\, 13,\, 17\}, \\
  S   &= \{101,\, \infty\}, \\
  L_T &= \Q(\sqrt{5},\, \sqrt{13},\, \sqrt{17},\, \sqrt{21},\, \sqrt{33}),
\end{align*}
with $[L_T : \Q] = 2^5 = 32$, and sets $K = L_T(i)$ as the CM field.

These values appear verbatim in the artifact:
\begin{itemize}[leftmargin=*]
  \item \texttt{T\_SET = (3, 5, 7, 11, 13, 17)}
  \item \texttt{S\_SET\_SPLIT = (101,)} (the infinite place is implicit)
  \item \texttt{L\_T\_GENERATOR\_RADICANDS = (5, 13, 17, 21, 33)}
  \item \texttt{SAWIN\_R\_PRODUCT = 2 \(\cdot\) 3 \(\cdot\) 5 \(\cdot\) 7 \(\cdot\) 11 \(\cdot\) 13 \(\cdot\) 17 = 510510}
\end{itemize}
and the test suite pins each constant against the published source.

\subsection{Splitting of 101 in $L_T$}

A rational prime $p$ splits completely in
$\Q(\sqrt{d_1}) \cdot \Q(\sqrt{d_2}) \cdots \Q(\sqrt{d_s})$ if and only
if it splits completely in each factor $\Q(\sqrt{d_i})$ (which, for
linearly disjoint quadratic extensions, is the definition of
completely splitting in the compositum). For odd
$p \neq d_i$, splitting in $\Q(\sqrt{d_i})$ is equivalent to
$\left(\frac{d_i}{p}\right) = 1$.

\begin{verification}\label{ver:101_splits}
The artifact evaluates the Jacobi symbol
$\left(\frac{d_i}{101}\right)$ for each
$d_i \in \{5, 13, 17, 21, 33\}$ and finds all five symbols equal to
$+1$. Therefore $101$ splits completely in $L_T$.
\textit{Source: \tname{src/erdos_ant/sawin_multiquadratic.py},
function \tname{splits_completely_in_L_T}.}
Pinned in
\tname{tests/test_sawin_multiquadratic.py::test_101_splits_completely_in_each_Q_sqrt_d}.
\end{verification}

\subsection{Galois-rank counts and Golod--Shafarevich admissibility}

Let $G = \Gal(\widetilde{\Q}_T / \Q)$ be the Galois group of the maximal
pro-$2$ extension of $\Q$ unramified outside $T$. Because $T$ contains a
prime that is $3 \pmod 4$ (e.g.\ $3$, $7$, or $11$), the $2$-rank of the
maximal elementary abelian unramified quotient is $|T| - 1 = 5$,
giving
\[
  d(G_T^\infty) = |T| - 1 = 5.
\]
Adjoining $S$ adds at most $|S| - 1 = 1$ Frobenius relations
(Sawin's pro-$2$ analogue of the Hajir--Maire--Ramakrishna cut), so
\[
  r(G_T^S) \leq d(G_T^\infty) + |S| - 1 = 5 + 1 = 6.
\]
The Golod--Shafarevich admissibility condition for the tower to remain
infinite is $r \leq d^2 / 4$; here
\[
  r(G_T^S) \leq 6 \leq \frac{(|T| - 1)^2}{4} = \frac{25}{4} = 6.25,
\]
which holds strictly. By the cited Golod--Shafarevich theorem
\cite{GolodShafarevich1964}, this admissibility implies the tower is
infinite. The artifact verifies the admissibility inequality only; it
does \emph{not} itself prove or compute the infinitude --- the cited
theorem does.

\begin{verification}\label{ver:gs_admissible}
The artifact computes $d(G_T^\infty) = 5$, $r(G_T^S) \leq 6$, and the
Golod--Shafarevich threshold $25/4 = 6.25$, and confirms
$6 \leq 6.25$. Pinned in
\tname{test_d_G_infty_equals_T_minus_1_because_T_contains_3_mod_4},
\tname{test_r_bound_is_six_per_remarks_pdf}, and
\tname{test_golod_shafarevich_inequality_holds}.
\end{verification}

\subsection{Non-claims}

Phase~3 does \emph{not} construct the infinite tower itself. The
existence of the tower is a non-constructive consequence of the
Golod--Shafarevich theorem \cite{GolodShafarevich1964} together with
Sawin's pro-$2$ analogue of the Hajir--Maire cut
\cite{HajirMaire2001,Sawin2026Remarks}; no finite computation can
exhibit a closed-form description of $L_j$ for arbitrary $j$. Phase~3
checks every step of the construction that \emph{is} amenable to finite
computation, and stops at that boundary.


% ---------------------------------------------------------------------------
\section{Numerical reproduction of equation~(2.2)}
\label{sec:eq22}

% Section 5: Numerical reproduction of remarks PDF eq (2.2) -- CENTRAL

This is the central numerical content of the artifact. We reproduce
equation~(2.2) of \cite{Sawin2026Remarks} as a regression test against
the published source.

\subsection{The published expression}

Sawin et al.\ state immediately after Lemma~2.1
\cite[\S2.1, eq.~(2.2)]{Sawin2026Remarks}:

\begin{quote}\itshape
Explicitly, the exponent in (2.1) can be taken to be
\[
  1 + \frac{\log(u \pi / (36 v))}{\log(36 / \delta^2)}
   \approx 1 + 6.24 \cdot 10^{-38}
\]
where $v = r/2$, $\delta \geq 101^{-2 \lceil 18 r^3 / \pi \rceil}$,
and $u = \lceil 18 r^3 / \pi \rceil \cdot r^{-2}$, with
$r = 2 \cdot 3 \cdot 5 \cdot 7 \cdot 11 \cdot 13 \cdot 17$.
\end{quote}

Let $K = \lceil 18 r^3 / \pi \rceil$, so that $u = K / r^2$ and
$\delta = 101^{-2K}$. Direct substitution gives
\[
  \frac{u \pi}{36 v} = \frac{(K / r^2) \cdot \pi}{36 \cdot (r/2)}
                    = \frac{K \pi}{18 r^3},
\]
and the exponent excess above $1$ simplifies to
\begin{equation}
  \delta_{\text{excess}} = \frac{\log(K \pi / (18 r^3))}{\log 36 + 4 K \log 101}.
  \label{eq:simplified}
\end{equation}

\subsection{Why \texttt{float64} fails}

The numerator of \eqref{eq:simplified} is $\log(K \pi / (18 r^3))$.
By definition of the ceiling, $K \geq 18 r^3 / \pi$, so the ratio
$K \pi / (18 r^3)$ lies in the half-open interval
\[
  \left(1,\; 1 + \frac{\pi}{18 r^3}\right] .
\]
For $r = 510510$, $r^3 \approx 1.33 \cdot 10^{17}$, hence
$\pi / (18 r^3) \approx 1.31 \cdot 10^{-18}$.

The numerator is therefore $\log(1 + \varepsilon)$ for some
$\varepsilon \in (0, 1.31 \cdot 10^{-18}]$. IEEE-754 \texttt{float64}
has a 52-bit mantissa and machine epsilon $\approx 2.22 \cdot 10^{-16}$,
so it cannot distinguish $1 + 1.3 \cdot 10^{-18}$ from $1$ in the first
place. The naive evaluation rounds the ratio to exactly $1$, yields
$\log 1 = 0$, and reports $\delta_{\text{excess}} = 0$.

This is not a flaw in eq.~\eqref{eq:simplified}; it is a property of
finite-precision arithmetic. Any honest evaluation must use a precision
significantly above $1$ part in $10^{18}$.

\subsection{Evaluation in \texttt{mpmath} at 200-bit precision}

The artifact uses \texttt{mpmath} \cite{mpmath} with
\texttt{mp.prec = 200} (about 60 decimal digits) to evaluate the
quantities in eq.~\eqref{eq:simplified} exactly enough that the
$1 + 10^{-18}$ separation survives. The function
\texttt{evaluate\_sawin\_exponent\_bound} in
\texttt{erdos\_ant.sawin\_multiquadratic} returns a
\texttt{SawinDeltaEvaluation} dataclass containing $r$, $K$, $v$,
$\log u$, and $\delta_{\text{excess}}$.

\begin{verification}\label{ver:eq22_reproduction}
Calling \tname{evaluate_sawin_exponent_bound()}
\textit{(source: \tname{src/erdos_ant/sawin_multiquadratic.py},
function \tname{evaluate_sawin_exponent_bound})} returns
\[
  r = 510{,}510, \qquad
  K = \lceil 18 \cdot 510510^3 / \pi \rceil \approx 7.62 \cdot 10^{17},
\]
and
\[
  \delta_{\text{excess}}
  = 6.2391 \cdot 10^{-38}.
\]
The published value is $\approx 6.24 \cdot 10^{-38}$ (given to two
significant figures in \cite{Sawin2026Remarks}). The relative error
between the artifact's evaluation and the published value is
\[
  \frac{|6.2391 - 6.24|}{6.24} = 1.4 \cdot 10^{-4} = 0.014\%,
\]
fully consistent with two-significant-figure rounding in the source
(see also \tname{docs/DEVIATION_LOG.md} for the full deviation
ledger).
Pinned in
\tname{tests/test_sawin_multiquadratic.py::test_sawin_exponent_bound_matches_remarks_eq_2_2}.
The per-source-file SHA-256 manifest under which this number was
computed appears in
\tname{reports/verification_result.json} field
\tname{source_sha256_manifest}.
\end{verification}

\subsection{What the verification establishes (and does not)}

Verification~\ref{ver:eq22_reproduction} establishes that, given the
inputs $T, S, L_T, K$ exactly as written in
\cite[\S2.1]{Sawin2026Remarks}, the elementary closed-form expression
eq.~\eqref{eq:simplified} evaluates to the value the source paper
claims, to within published rounding.

This does \emph{not} establish:
\begin{itemize}[leftmargin=*]
  \item that the closed-form expression is itself correct (that
        depends on Lemma~2.1 and earlier algebra of \cite{Sawin2026Remarks},
        which we do not re-derive);
  \item that the underlying infinite Golod--Shafarevich tower exists
        (Phase~3 only verifies the finite admissibility condition);
  \item that $\delta_{\text{excess}} = 6.24 \cdot 10^{-38}$ is in any
        sense optimal --- it is one explicit lower bound, not the best
        possible.
\end{itemize}

What it does establish is that any future code change in the artifact
that drifts away from the published value will fail the regression
test at the $10^{-3}$ relative-tolerance band hard-coded in the test
(currently $0.014\%$, below the tolerance but close enough to catch
material drift).


% ---------------------------------------------------------------------------
\section{Limits and explicit non-claims}
\label{sec:limits}

% Section 6: Limits and explicit non-claims

This section is placed before reproducibility deliberately. The shortest
honest summary of this artifact is the list of things it does not do.

\subsection{This artifact does NOT contain a new proof}

Theorem~1.1 of \cite{OpenAI2026Proof,Sawin2026Remarks} is proved by
those authors. The present paper does not re-derive Lemma~2.1,
Lemma~2.2, Proposition~2.3, or any other intermediate result. Where the
source papers say ``by the Golod--Shafarevich theorem'' or ``by the
Hajir--Maire tower-cutting argument'', we likewise take those classical
results as given.

\subsection{This artifact does NOT construct the infinite tower}

The existence of $L_j$ with $[L_j : \Q] \to \infty$ and bounded root
discriminant is supplied by Golod--Shafarevich
\cite{GolodShafarevich1964} together with the Hajir--Maire splitting
technique \cite{HajirMaire2001} adapted to pro-$2$ in
\cite{Sawin2026Remarks}. This is a non-constructive existence theorem.
No finite computation can exhibit the entire infinite tower. Phase~3
of the artifact stops at the boundary of computable
admissibility.

\subsection{This artifact has NOT been peer-reviewed}

No external mathematician has reviewed this code base. Our confidence
in the artifact rests on three independent signals:

\begin{enumerate}[leftmargin=*]
  \item the 59-test unit suite passes on CI (Ubuntu + Windows
        $\times$ Python 3.11 / 3.12);
  \item the central numerical reproduction matches the published
        $\approx 6.24 \cdot 10^{-38}$ to $0.014\%$ relative error
        (\S\ref{sec:eq22});
  \item the source code is small (\(\approx\) 1\,000 LOC across four
        modules) and intended to be readable by anyone who has read
        \cite{Sawin2026Remarks} \S2.
\end{enumerate}

None of these substitutes for mathematician review. The artifact is
offered as input to such review, not as a substitute.

\subsection{Reported $0.014$ and $0.0318$ are NOT in either formal PDF}

Secondary coverage (e.g.\ \cite{KalaiBlog2026}) reports numerical
values $\delta \approx 0.014$ and $\delta \approx 0.0318$ in the
informal contexts $n^{1.014}$ and $n^{1.0318}$. We verified by direct
text extraction with \texttt{pypdf} that neither of these literal
numerical values appears in either
\cite{OpenAI2026Proof} or \cite{Sawin2026Remarks}. The only rigorous
numerical lower bound that appears verbatim in either formal PDF is
$\delta_{\text{excess}} \approx 6.24 \cdot 10^{-38}$ from eq.~(2.2) of
\cite{Sawin2026Remarks}.

The artifact's \texttt{SOURCE\_CLAIMS} registry encodes this
three-way provenance explicitly:

\begin{itemize}[leftmargin=*]
  \item \texttt{openai\_index\_reported} ($\delta = 0.014$,
        \texttt{verified\_numerical\_in\_formal\_pdf = False}),
  \item \texttt{sawin\_rigorous\_lower\_bound}
        ($\delta = 6.24 \cdot 10^{-38}$,
        \texttt{verified\_numerical\_in\_formal\_pdf = True}),
  \item \texttt{sawin\_refinement\_reported} ($\delta = 0.0318$,
        \texttt{verified\_numerical\_in\_formal\_pdf = False}).
\end{itemize}

\subsection{This artifact does NOT improve the published bound}

We made no attempt to optimise the choice of $T$, $S$, $L_T$ or the
inner parameter $k$ to obtain a larger $\delta_{\text{excess}}$. The
explicit choice of \cite[\S2.1]{Sawin2026Remarks} is what we encode and
reproduce; the optimisation problem of finding $(T, S, k)$ with larger
admissible $\delta$ remains open and is out of scope.


% ---------------------------------------------------------------------------
\section{Reproducibility}
\label{sec:repro}

% Section 7: Reproducibility

\subsection{One-line reproduction}

\begin{quote}\ttfamily
git clone https://github.com/Flamehaven-Labs/openai-erdos-eq22-reproduction\\
cd openai-erdos-eq22-reproduction\\
pip install -e ".[dev]"\\
python -m erdos\_ant.verify
\end{quote}

\noindent
Expected terminal output:

\begin{quote}\ttfamily
Verdict: PASS\\
Checks: 21/21 passed\\
eq (2.2) exponent excess: 6.2391e-38 (published \textasciitilde 6.24e-38, rel.err 0.0001)\\
(frozen-report mode: tracked evidence not written; pass --write-evidence to refresh)
\end{quote}

\noindent
By default the command does not write to disk; the working tree stays
clean. Release maintainers pass \texttt{--write-evidence} to refresh the
tracked files
\path{reports/verification_result.json} and
\path{reports/verification_report.md}.

The line \texttt{Checks: 21/21 passed} refers to the verifier's
top-level checklist. One of those checks is a subprocess execution of
the $60$-test pytest suite with external pytest plugin autoloading
disabled; running \texttt{pytest -q} directly executes the same unit
tests outside the verifier wrapper.

\subsection{Pinned environment}

\begin{tabularx}{\textwidth}{@{}p{0.30\textwidth}X@{}}
\toprule
Item & Value \\
\midrule
Repository & \url{https://github.com/Flamehaven-Labs/openai-erdos-eq22-reproduction} \\
Version (paper source) & v0.2.5 \\
Python & 3.11 and 3.12 (CI-verified) \\
numpy & $\geq 1.26$ \\
mpmath & $\geq 1.3$ \\
Test count & $60$ unit tests $+$ $21$ verifier checks \\
CI matrix & Ubuntu $\times$ Windows $\times$ Python 3.11 / 3.12 \\
Licence & MIT \\
Per-release DOIs & \texttt{CITATION.cff} and \href{https://github.com/Flamehaven-Labs/openai-erdos-eq22-reproduction/releases}{GitHub Releases} \\
\bottomrule
\end{tabularx}

\subsection{Supplementary local checks}

In addition to the unit suite and verifier, the release process keeps
author-operated local source-level checks for claim framing,
traceability, and repository hygiene. Those checks are recorded as
engineering evidence only. They are not independent review, are not
mathematical peer review, and are not used as a premise for the
correctness of the eq.~(2.2) reproduction. See \S\ref{sec:limits} for
the bound on what local release checks can and cannot attest to.

\subsection{Citation}

A \texttt{CITATION.cff} file in the repository root provides citation
metadata in the Citation File Format v1.2 standard, recognised by
Zenodo, GitHub, and most citation managers. The recommended citation
for a specific tagged release is the Zenodo DOI associated with that
tag.


% ---------------------------------------------------------------------------
\clearpage
\section*{Acknowledgements}

% Section 8: Acknowledgements

We thank the OpenAI reasoning model whose May 2026 output prompted
this reproduction effort, and the authors of the companion remarks
paper \cite{Sawin2026Remarks} --- Noga Alon, Thomas F.~Bloom,
W.~T.~Gowers, Daniel Litt, Will Sawin, Arul Shankar, Jacob Tsimerman,
Victor Wang, and Melanie Matchett Wood --- for the explicit choice of
parameters in their \S2 that made a finite numerical reproduction
tractable in the first place.

We thank Farshid Hajir, Christian Maire, and Ravi Ramakrishna for the
class-field tower-cutting technique \cite{HajirMaire2001} that is
invoked, but not re-derived, in this artifact's Phase~3; and David
A.~Cox, whose textbook treatment of genus theory underlies Phase~2.

This artifact was developed and released by the author independently.
The supplementary local checks listed in \S\ref{sec:repro} were run as
part of the release process; their outputs are treated as recorded
engineering evidence, not as independent mathematical endorsement.


% ---------------------------------------------------------------------------
\clearpage
\printbibliography

\end{document}
